\section{考察}
\label{sec:discussion}

\emph{実用上の指針．}
感情ごとに少数のラベルしかない新しいコーパスに対して，結果は次の4つの手順を示唆する．
すなわち，ソースとターゲットの特徴を別々に中心化してラベルを使わずに判断基準による損失の大部分を取り除き（\ref{sec:part1-fix}節），すべてのソースから，切片には緩いアンカーを，重みには強いアンカーをかけて出発し（\ref{sec:part2-start}節），どの発話にラベルを付けるかをラベルを使わずに選んでもよい出発点の向きが直ることは期待せず（\ref{sec:part2-limits}節），ターゲットの発話とともに他のモダリティが記録されているならば，それがラベル付けのときだけであっても，ソースで学習し，それらのモダリティを見る先生にプール全体のラベルを付けさせ，配備する分類器は単一モダリティのままにする（\ref{sec:comparison}節）．

\emph{テキストのふるまいが異なる理由．}
テキストの特徴は判断基準による損失が小さく，向きによる損失が大きい（表\ref{tab:part1}）ので，少数のラベルには直すべき判断基準のずれがほとんど無く，利得のうち判断基準の割合は最大で\NUM{p1-e29-critshare-k1-t}である．
利得のうち判断基準から生じる割合は，損失のうち判断基準が引き起こす割合に従う．
それでもテキストの生徒は，テキストのラベルに欠けている向きの情報を補う融合の先生から利得を得る．

\emph{エントロピー最小化とクラスタの仮定．}
エントロピー最小化は，ラベルのないデータでクラスが分かれたクラスタをなすことに依拠し，この性質はAVTとATの構成のIEMOCAPの即興のセッションでは成り立つが，MELD，CMU-MOSEI，およびVTとAVの構成の同じターゲットでは成り立たない（\ref{sec:comparison-fixed}節）．
別のモダリティからの弱いラベルはモダリティ間の相補性に依拠しており，その利得のばらつきはずっと小さい．
$K=0$と$K=5$のターゲットとセルの\NUM{st-b-tgt-n}の組み合わせのうち，提案法の拡張版が出発点を下回るのは\NUM{st-b-tgtneg-n}で，最も悪い差は\NUM{st-b-tgtneg-worst}ポイントである．
